\chapter{Autoantibody Testing}

\section*{What Is This Test?}
Autoantibody testing detects antibodies directed against the body's own cells, tissues, or proteins. These tests support evaluation of autoimmune disease but usually cannot establish a diagnosis by themselves.

\section*{Alternative Names}
\begin{itemize}
  \item Autoimmune Antibody Testing
  \item Autoimmune Serology
  \item Autoantibody Panel
  \item Rheumatologic Serology
\end{itemize}

\section*{Principle}
The laboratory uses immunoassays, indirect immunofluorescence, immunoblotting, or related methods to detect a specific antibody. Common examples include ANA, anti-dsDNA, ENA antibodies, ANCA, rheumatoid factor, anti-CCP, antiphospholipid antibodies, and organ-specific antibodies.

\section*{Why This Test Is Ordered}
Testing is selected for symptoms or findings that suggest a particular autoimmune condition, such as inflammatory arthritis, connective-tissue disease, vasculitis, autoimmune liver disease, thyroid disease, or autoimmune cytopenia. Broad panels without a compatible clinical question can produce confusing false-positive results.

\section*{Primary vs Secondary Test}
The appropriate antibody assay is a secondary diagnostic test guided by the history, examination, and initial laboratory findings. ANA is a screening test for some connective-tissue diseases, but a positive ANA is not equivalent to an autoimmune diagnosis.

\section*{How This Test Is Performed}
A venous blood sample is collected and tested using the method selected by the laboratory. Some tests report a concentration or index, while ANA by immunofluorescence may also report a titer and staining pattern.

\section*{What Preparation Is Needed}
Fasting is usually not required. Tell the clinician about autoimmune medicines, infections, vaccinations, pregnancy, and relevant personal or family history.

\section*{Understanding Results}
Interpretation depends on antibody specificity, level, pattern, laboratory cutoff, symptoms, and organ findings. Low-level or isolated positivity can occur in healthy people or after infection. A negative result reduces the likelihood of some diseases but does not exclude all autoimmune disease.

\section*{Follow-up Tests}
Follow-up may include CBC, urinalysis, kidney and liver tests, complement levels, inflammatory markers, disease-specific antibodies, imaging, or tissue biopsy. Urgent assessment is needed for suspected organ-threatening disease.

\section*{References}
\begin{enumerate}
  \item MedlinePlus. Autoantibody Testing.
  \item American College of Rheumatology. Understanding Your ANA Test.
  \item International Consensus on ANA Patterns.
\end{enumerate}

\clearpage
\section*{Understanding Results and Follow-up}
The laboratory report should identify the specimen, method, units, reference interval or decision limit, and whether the result is positive, negative, abnormal, or inconclusive. Reference intervals vary with age, sex, pregnancy, specimen type, assay, and laboratory; a result outside a range is not automatically a diagnosis, and a result within a range does not always exclude disease. Discuss unexpected results with the ordering clinician, who can decide whether repeat or confirmatory testing, treatment, or referral is appropriate. Do not change medicines or delay urgent care based on an isolated result.


\section*{Regulatory Status and Authorization Date}
This chapter describes a test category, not a specific commercial product. FDA, CDSCO, EU IVDR, and MHRA approval or registration dates therefore cannot be assigned without the assay manufacturer, product name, instrument, and intended use. For laboratory-developed tests, an individual agency approval date may not exist. See the Preface for the interpretation of regulatory status.

\clearpage
